% !TeX root = ../../main.tex
% chapters/5_mobiwac/08_conclusion.tex -- one section of Chapter 5 (MobiWac 2026, submitted, under review).
% Split out of chapters/5_mobiwac.tex on 2026-07-28 (author request: the paper
% chapters were single long files, while the originals under articles/[mobiwac]/src/sections/
% are one file per section). The split is PURELY MECHANICAL: content was cut at
% \section boundaries and not edited. The chapter file now \inputs these in order.
% Editing rules are unchanged and still governed by the chapter's status above:
% see NORTH_STAR section 4 before changing any sentence here.
\section{Conclusion}
\label{sec:mobiwac:conclusion}

We tested whether one model can predict both the category and region of a user's next visit. A check-in-level
representation, one vector per visit rather than one per place, makes the next-category task far more learnable than the standard
place-level representation does. On that representation, a single multitask model outperforms a dedicated category model on every dataset. On region, it outperforms the dedicated model on four of the six datasets and remains non-inferior (TOST, $\pm2$~pp) on the other two. The gains at Texas and California exceed the two-point margin; those at Florida and Istanbul are statistically supported but smaller. Across the U.S. states, the region gain rises with region count. The joint model also remains above every external baseline in Table~\ref{tab:mobiwac:results} on all six datasets, by at least 4 Acc@10 points over the strongest region reference (HMT-GRN, STAN, or ReHDM; the first two on our folds, ReHDM under its own protocol) and by at least 33 macro-F1 points over POI-RGNN on category. One model with a shared semantic context and a private spatial path anticipates both what a user will do next and where.
